\subsection*{CU-49: Eliminar a un amigo}
\addcontentsline{toc}{subsection}{CU-49: Eliminar a un amigo}
\label{cu-49}

\begin{fichacu}{CU-49}{Eliminar a un amigo}
  \crudFila{Responsable}{Ángel Gabriel Aguilar Hernández}
  \crudFila{Actor principal}{Jugador}
  \crudFila{Actores de sistema}{Servidor de partidas, Base de datos}
  \crudFila{Prototipos}{2a, 2b, 5b}
  \crudFila{Objetivo}{Permitir al Jugador quitar a otro jugador de su lista de amigos sin bloquearlo y sin avisarle, de modo que ambos dejen de figurar como amigos pero puedan volver a serlo.}
  \crudFila{Disparador}{El Jugador selecciona ``eliminar amigo'' en el menú de un amigo en la \texttt{GUI\_\allowbreak{}Friends} o en su tarjeta de perfil.}
  \textbf{Precondiciones} & \cuCelda{\begin{cureglas}
      \item \textbf{\mbox{PRE-01}} El Jugador tiene una \textit{Sesion} abierta y su token es válido.
      \item \textbf{\mbox{PRE-02}} Existe una \textit{Amistad} entre el Jugador y el otro jugador.
      \item \textbf{\mbox{PRE-03}} El Servidor de partidas está en ejecución y tiene conexión disponible con la Base de datos.
    \end{cureglas}} \\ \hline
  \textbf{Flujo normal} & \cuCelda{\begin{cuenum}[start=1]
      \item El Jugador abre el menú de un amigo y selecciona ``eliminar amigo''.
      \item El SJM despliega la \texttt{GUI\_\allowbreak{}MessageConfirm} indicando que dejarán de ser amigos para ambos, que el otro no recibirá ningún aviso y que podrán volver a enviarse solicitud, con las opciones ``Eliminar'', ``Eliminar y bloquear'' y ``Cancelar''. \textit{(Ver \mbox{FA-01}, \mbox{FA-02})}
      \item El Jugador selecciona ``Eliminar''.
      \item El SJM envía el mensaje \texttt{REMOVE\_\allowbreak{}FRIEND\_\allowbreak{}REQUEST} con el \texttt{id\_\allowbreak{}usuario} del amigo y el token. \textit{(Ver \mbox{EX-02}, \mbox{EX-03}, \mbox{EX-04})}
      \item El Servidor de partidas verifica que el token corresponda a una \textit{Sesion} abierta. \textit{(Ver \mbox{FA-03})}
      \item El Servidor de partidas ordena el par de identificadores y busca la \textit{Amistad} con el menor en \texttt{id\_\allowbreak{}usuario\_\allowbreak{}a} y el mayor en \texttt{id\_\allowbreak{}usuario\_\allowbreak{}b} (\mbox{RN-02}). \textit{(Ver \mbox{FA-04})}
      \item El Servidor de partidas elimina la \textit{Amistad}. \textit{(Ver \mbox{EX-01}, \mbox{FA-05}, \mbox{FA-06})}
      \item El Servidor de partidas \textbf{no} notifica al otro jugador (\mbox{RN-03}).
    \end{cuenum}} \\
   & \cuCelda{\begin{cuenum}[start=9]
      \item El Servidor de partidas responde con \texttt{REMOVE\_\allowbreak{}FRIEND\_\allowbreak{}OK}.
      \item El SJM retira al jugador de la lista de amigos y del bloque de quién está jugando ahora.
      \item Fin del caso de uso.
    \end{cuenum}} \\ \hline
  \textbf{Flujos alternos} & \cuCelda{\textbf{\mbox{FA-01}: El Jugador cancela}\par
    \begin{cuenum}[start=1]
      \item El Jugador selecciona ``Cancelar''.
      \item El SJM cierra la confirmación sin enviar nada.
      \item Fin del caso de uso.
    \end{cuenum}} \\
   & \cuCelda{\textbf{\mbox{FA-02}: El Jugador elimina y bloquea}\par
    \begin{cuenum}[start=1]
      \item El Jugador selecciona ``Eliminar y bloquear''.
      \item El flujo continúa en el paso 4 del \mbox{CU-33} Silenciar o bloquear a un jugador, que elimina la amistad dentro de la misma transacción del bloqueo.
      \item Fin del caso de uso.
    \end{cuenum}} \\
   & \cuCelda{\textbf{\mbox{FA-03}: La sesión ya no es válida}\par
    \begin{cuenum}[start=1]
      \item El Servidor de partidas responde con \texttt{SESSION\_\allowbreak{}INVALID}.
      \item El SJM muestra la \texttt{GUI\_\allowbreak{}MessageWarning} indicando que la sesión terminó y despliega la \texttt{GUI\_\allowbreak{}Login}.
      \item Fin del caso de uso.
    \end{cuenum}} \\
   & \cuCelda{\textbf{\mbox{FA-04}: Ya no eran amigos}\par
    \begin{cuenum}[start=1]
      \item El Servidor de partidas no encuentra la \textit{Amistad}: el otro jugador la eliminó antes, lo bloqueó o su cuenta se dio de baja.
      \item El Servidor de partidas responde con \texttt{REMOVE\_\allowbreak{}FRIEND\_\allowbreak{}OK}, porque el estado que el Jugador pedía ya se cumple.
      \item El flujo continúa en el paso 10 del FN.
    \end{cuenum}} \\
   & \cuCelda{\textbf{\mbox{FA-05}: El amigo está en la misma sala o partida que el Jugador}\par
    \begin{cuenum}[start=1]
      \item El Servidor de partidas elimina la \textit{Amistad} con normalidad.
      \item La sala o la partida en curso no se ven afectadas: la amistad no es condición para jugar juntos (\mbox{RN-04}).
      \item El flujo continúa en el paso 9 del FN.
    \end{cuenum}} \\
   & \cuCelda{\textbf{\mbox{FA-06}: Hay una invitación pendiente entre ambos}\par
    \begin{cuenum}[start=1]
      \item El Servidor de partidas elimina la \textit{Amistad} y conserva la \textit{Invitacion} pendiente, que expirará sola en menos de un minuto: invitar no exige ser amigos (\mbox{RN-02} del \mbox{CU-32}).
      \item El flujo continúa en el paso 9 del FN.
    \end{cuenum}} \\ \hline
  \textbf{Excepciones} & \cuCelda{\textbf{\mbox{EX-01}: Fallo al escribir en la base de datos}\par
    \begin{cuenum}[start=1]
      \item El Servidor de partidas registra la excepción con un identificador de incidencia y responde con \texttt{SERVER\_\allowbreak{}ERROR}.
      \item El SJM muestra la \texttt{GUI\_\allowbreak{}MessageError} con el identificador y conserva al amigo en la lista.
      \item Fin del caso de uso.
    \end{cuenum}} \\
   & \cuCelda{\textbf{\mbox{EX-02}: Se agota el tiempo de espera de la respuesta}\par
    \begin{cuenum}[start=1]
      \item El SJM no reenvía la solicitud y recarga la lista de amigos, que dirá si se aplicó.
      \item Fin del caso de uso.
    \end{cuenum}} \\
   & \cuCelda{\textbf{\mbox{EX-03}: La conexión se interrumpe durante el intercambio}\par
    \begin{cuenum}[start=1]
      \item El SJM muestra la barra de conexión perdida y recarga la lista al recuperarla.
      \item Fin del caso de uso.
    \end{cuenum}} \\
   & \cuCelda{\textbf{\mbox{EX-04}: El mensaje recibido está malformado}\par
    \begin{cuenum}[start=1]
      \item El SJM descarta el mensaje, cierra la conexión y registra en la bitácora local el contenido truncado.
      \item El SJM muestra la \texttt{GUI\_\allowbreak{}MessageError} indicando que la respuesta no pudo interpretarse.
      \item Fin del caso de uso.
    \end{cuenum}} \\ \hline
  \textbf{Postcondiciones} & \cuCelda{\begin{cureglas}
      \item \textbf{\mbox{POST-01}} Si el caso termina por el FN, no existe ninguna \textit{Amistad} entre ambos jugadores.
      \item \textbf{\mbox{POST-02}} Ninguno de los dos aparece en la lista de amigos del otro.
      \item \textbf{\mbox{POST-03}} No se creó ningún \textit{Bloqueo} ni \textit{Silencio}, salvo en el \mbox{FA-02}.
      \item \textbf{\mbox{POST-04}} Las \textit{Participacion} de las partidas que jugaron juntos, y por tanto el marcador entre ambos, se conservan.
      \item \textbf{\mbox{POST-05}} Ambos pueden volver a enviarse una solicitud de amistad con el \mbox{CU-30}.
    \end{cureglas}} \\ \hline
  \textbf{Reglas de negocio} & \cuCelda{\begin{cureglas}
      \item \textbf{\mbox{RN-01}} Eliminar a un amigo es recíproco: la amistad desaparece para los dos.
      \item \textbf{\mbox{RN-02}} La \textit{Amistad} se busca por el par ordenado, igual que se guardó en el \mbox{CU-31}, así que da igual cuál de los dos la elimine.
      \item \textbf{\mbox{RN-03}} El otro jugador no recibe ningún aviso.
      \item \textbf{\mbox{RN-04}} La amistad no condiciona jugar juntos: eliminarla no afecta a salas, partidas ni invitaciones en curso.
      \item \textbf{\mbox{RN-05}} Eliminar a un amigo no lo bloquea. Para cortar el contacto está el \mbox{CU-33}.
      \item \textbf{\mbox{RN-06}} La \textit{Amistad} se elimina físicamente: describe una relación presente, no un hecho histórico.
    \end{cureglas}} \\ \hline
  \textbf{Observación} & \cuCelda{\textit{Sobre por qué este caso se añadió.}\par
    El catálogo original tenía casos para crear una amistad pero ninguno para terminarla; la única forma de dejar de ser amigos era bloquear, que corta mucho más de lo que el jugador quiere. El análisis CRUD lo puso en evidencia: \textit{Amistad} tenía creación y consulta pero no eliminación. Es el caso más simple del documento, y aun así evita que la tabla solo pueda crecer.} \\ \hline
  \textbf{Datos que toca} & \cuCelda{\begin{cudatos}
      \item \textit{Sesion}: lee por token \textit{(paso 5)}
      \item \textit{Amistad}: lee por par ordenado; elimina \textit{(pasos 6, 7)}
    \end{cudatos}} \\ \hline
\end{fichacu}
\clearpage
